% ===========================================================================
\section{Introdução}
% ===========================================================================

A multiplicação de inteiros é uma operação com pouca estrutura visível. A adição
na reta real tem ordem, distância e simetria; o produto em $\N$ tem apenas a
fatoração. O logaritmo é a ponte entre as duas: ele transforma produto em soma,
e é a única (a menos de escala) com essa propriedade. Este artigo percorre essa
ponte até o fim, e o que se encontra do outro lado é uma imagem — literalmente,
uma imagem em frequência — da aritmética elementar.

A pergunta de partida é simples de enunciar. Todo inteiro $m \ge 2$ tem uma
frequência natural,
\[
  \nfreq{m} \;=\; \log m ,
\]
e a fatoração única diz que $m$ é composto exatamente quando $\nfreq{m}$ pode ser
escrito como soma de frequências de inteiros menores. Ou seja: \emph{a
primalidade é a irredutibilidade de um ponto num semigrupo aditivo}. Isso, por si
só, é apenas uma reformulação. A reformulação torna-se interessante quando se
pergunta pela \emph{distância} — quão longe $\nfreq{m}$ está do semigrupo gerado
por um conjunto de primos — porque aí a aritmética ganha uma geometria, e a
geometria pode ser desenhada.

O artigo constrói três coisas.

\paragraph{O dicionário.} Na Seção~\ref{sec:dicionario} formalizamos a tradução
entre a estrutura multiplicativa de $\N$ e a estrutura aditiva da reta. O ponto
central é que ela é um isomorfismo de semigrupos sobre os números gerados por um
conjunto de primos, de modo que perguntas multiplicativas (``$m$ é divisível por
$p$?'', ``$m$ é primo?'') se tornam perguntas aditivas (``$\nfreq{m}$ está no
reticulado?'', ``$\nfreq{m}$ é gerador?'').

\paragraph{O crivo como partição.} Na Seção~\ref{sec:blocos} mostramos que a
partição de $\N$ em blocos binários
\[
  A_k \;=\; [\,2^k,\; 2^{k+1}) , \qquad k \ge 0 ,
\]
é a partição adaptada ao dicionário. O fato decisivo é elementar: \emph{todo
composto de $A_k$ tem um fator primo menor que $2^k$} (Lema~\ref{lema:bloco}). A
consequência é que cada bloco pode ser classificado usando apenas o bloco
anterior como base — e que essa base, a \emph{semente} $\Seed_k$, é exatamente a
que o dicionário exige. O crivo de Eratóstenes aparece então como um caso
particular de uma recorrência sobre blocos, sem nada de espectral.

\paragraph{A geometria exata.} A Seção~\ref{sec:geometria} é o núcleo matemático.
Definimos a distância relativa $\rB_{\Seed}(m)$ de $\nfreq{m}$ ao reticulado
gerado pelos logaritmos de $\Seed$, e determinamos seu valor \emph{exato} para os
primos de um bloco:
\begin{equation}\label{eq:rhoexato}
  \rB(q) \;=\; \frac{\log\bigl(1+1/q\bigr)}{\log q}
  \qquad \text{para todo primo } q \in A_k,\ k\ge 2 .
\end{equation}
Não é uma cota: é igualdade. E a razão é uma única observação, quase banal — os
dois vizinhos inteiros de um primo ímpar são compostos, logo ambos pertencem ao
reticulado, de modo que o ponto do reticulado mais próximo de $\nfreq{q}$ é sempre
$\log(q+1)$. O reticulado gerado pelos primos de um bloco é, na escala $1/q$,
indistinguível do conjunto de \emph{todos} os logaritmos de inteiros. Desse valor
exato saem a cota inferior uniforme
$\rB(q) \ge \bigl(2^{k+1}(k+1)\log 2\bigr)^{-1}$ para $q \in A_k$, o mínimo
$\rB_{\min}(k) = \log(1+1/q_k)/\log q_k$ (atingido no maior primo $q_k$ do bloco)
e o assintótico $\rB_{\min}(k) \sim 1/(q_k \log q_k)$.

Com isso o dicionário fica completo, e o retrato que ele permite desenhar é o das
Seções~\ref{sec:desenho} a \ref{sec:assinatura}. Três retratos, na verdade.

\paragraph{Primeiro retrato: o desenho do crivo.} Seção~\ref{sec:desenho}. Para
cada inteiro $m$ considere a função
\[
  S_m(t) \;=\; -\tfrac12 \log\!\Bigl(1 - 2\,m^{-1/2}\cos\bigl(t\nfreq{m}\bigr)
  + m^{-1}\Bigr) ,
\]
que é o logaritmo do módulo do fator $\bigl(1-m^{-s}\bigr)^{-1}$ avaliado em
$s=\tfrac12+it$. Ela oscila com frequência $\nfreq{m}/(2\pi)$ e amplitude
$\approx m^{-1/2}$. Somando sobre um bloco e subtraindo a parte devida aos
compostos, o resíduo é a soma dos $S_q$ sobre os primos do bloco — e a
transformada de Fourier desse resíduo tem exatamente um pico por primo, em
posição $\nfreq{q}/(2\pi)$ (Figura~\ref{fig:desenho}). A subtração é uma operação
aritmética simples e a imagem é fiel; o que a figura mostra é o bloco sem os
compostos, e nada mais.

\paragraph{Segundo retrato: a lei de resolução.} Seção~\ref{sec:resolucao}. Um
desenho tem custo. Para separar dois primos consecutivos $q_1<q_2$ na imagem é
necessário observar o sinal até $t_{\max} > 2\pi/\log(q_2/q_1)$ (critério de
Rayleigh). O par que impõe o maior requisito é o de menor \emph{gap relativo} —
em geral um par de gêmeos — e daí resulta que o custo do desenho de um bloco é
$\Theta(2^n/n)$, com cotas explícitas (Teorema~\ref{teo:custo}). No nível de um
primo isolado, o requisito é $\Theta(q)$: ver a irredutibilidade de $q$ custa uma
janela de tamanho proporcional a $q$. Essa é uma propriedade da representação, e
responde a uma pergunta natural de quem olha as figuras: quanto é preciso
``aumentar o zoom''.

\paragraph{Terceiro retrato: a contagem.} Seção~\ref{sec:contagem}. Há um caso em
que a representação espectral faz algo que a aritmética faz de outro modo, e o faz
com exatidão absoluta. Remova de cada inteiro de um lote o seu núcleo, isto é, a
parte formada por um conjunto fixo $P$ de primos pequenos. Se $k$ candidatos do
lote têm o mesmo núcleo $r$, então os $k$ termos resultantes são literalmente a
\emph{mesma} função de $t$, e somam $k$ vezes. A amplitude em $\nfreq r$ é
exatamente $k$ vezes a de um único termo, para qualquer janela e qualquer
resolução (Proposição~\ref{prop:ganho}). O espectro passa a ser um contador:
a altura de um pico é o número de candidatos que compartilham aquele núcleo.

\paragraph{E um quarto, menor.} Seção~\ref{sec:assinatura}. Para um par
$(p,p+2)$, a presença do tom em $\nfreq{p+2}$ é detectável com contraste que cresce
como $T/p$; pares gêmeos e pares mistos se separam, e o fundo decai como
$1/T$. É a versão local do segundo retrato.

\medskip

\paragraph{Escopo.} O objetivo do artigo é a \emph{representação}, não a
eficiência: medimos o custo de ver, e não propomos substituir a aritmética por
transformadas. A aritmética decide a primalidade de $m$ em $O(\sqrt m)$ operações;
a imagem espectral exige uma janela de tamanho $\Theta(m)$. O valor do dicionário
está em outro lugar: ele torna visíveis relações que a notação multiplicativa
esconde, e dá a cada uma delas uma lei quantitativa. As figuras deste artigo foram
feitas para serem olhadas.

\paragraph{Convenções.} $\N=\{1,2,3,\dots\}$; $p$ e $q$ sempre denotam primos;
$\log$ é o logaritmo natural; $v_p(m)$ é o expoente de $p$ em $m$. Escrevemos
$[a,b)=\{m\in\N: a\le m<b\}$ e $\#X$ para a cardinalidade de um conjunto finito
$X$. O Teorema Fundamental da Aritmética é usado livremente, e apenas ele.
